\section{Progettazione Ad Alto Livello}
\label{sec:ProgettazioneAdAltoLivello}

\subsection{Front-end}
\label{sec:Front-end}
\subsubsection{StandardPage}
\label{sec:StandardPage}
\paragraph{Descrizione}
La StandardPage è la pagina che funge da struttura del layout che ogni pagina riceve come parametro con il nome del componente BasePageComponent.
\paragraph{SottoComponenti}
StandardPage è composto dai seguenti sotto-componenti:
\begin{itemize}
    \item \textbf{BasePageComponent}: Il componente con dichiarati i metodi e parametri necessari alla creazione di una StandardPage tra i quali:
    \begin{itemize}
        \item \textbf{setHeader()}: Metodo che va a chiamare i componenti dediti alla costruzione dell'header.
        \item \textbf{setFooter()}: Metodo che va a chiamare i componenti dediti alla costruzione dell'footer.
        \item \textbf{setMessages()}: Metodo che va a chiamare i componenti dediti alla costruzione di messaggi di sistema tra qui possibili messaggi a comparsa, pop-up e messaggi d'errore standardizzabili.
        \item \textbf{setLoadingStatus()}: Metodo che va a chiamare i componenti dediti al feedback del caricamento a livello grafico, ogni pagina dunque sarà capace di avviare un caricamento.
    \end{itemize}
    \item \textbf{ObserverComponent}: Componente dotato del metodo sovracscrivibile Update().
    \item \textbf{SubjectComponent}: Componente dotato del metodo sovracscrivibile Subscribe().
    \item \textbf{NavigationLinksComponent}: Componente dotato del metodo Update() che eredita da ObserverComponent, tale componente si occupa di aggiornare i link presenti sull header a seconda della pagina richiesta.
    \item \textbf{HeaderComponent}: Componente dotato del metodo Subscribe() che eredita da SubjectComponent, tale componente riceve le modifiche richieste dal NavigationLinksComponent per aggiornare l'header.
    \item \textbf{FooterComponent}: Componente che instanzia il footer sul ngOnInit().
    \item \textbf{LoadingComponent}: Il componente con dichiarati i metodi necessari all'avvio del caricamento grafico:
    \begin{itemize}
        \item \textbf{loadingStatus}: Metodo booleano che indica la presenza o meno del caricamento.
        \item \textbf{ngOnInit()}: Metodo che si occupa si settare il parametro loadingStatus a False.
        \item \textbf{setStatus()}: Metodo che imposta il valore del parametro loadingStatus.
    \end{itemize}
    \item \textbf{MessageHandlerComponent}: Componente dotato del metodo Update() che eredita da ObserverComponent, tale componente si occupa di creare il tipo e contenuto del messaggio.
    \item \textbf{MessageComponent}: Componente dotato del metodo Subscribe() che eredita da SubjectComponent, tale componente riceve le impostazioni del messaggio richieste dal MessageHandlerComponent per aggiornare l'apparizione del messaggio.
\end{itemize}
